% Transcription — fonds Alexandre Grothendieck, shelfmark 57, batch 1 (pages 1-20).
% Established from the facsimile archives/batches/57/batch-01.pdf, per the skill
% transcribe-grothendieck.
%
% Pass: Opus 5.5 (claude-opus-5-5), 2026-10-03 — first pass, unchecked against the pages by a human.
%
% Sixteen of the twenty pages are transcribed. Page 1 is blank and absent;
% page 13 is a cover sheet in his hand (« Carré, cube, rigidité… »), which
% becomes the heading of the run that follows and gets no \page{}.
%
%   Pages 2-12 — an English typescript, « Local Picard schemes », in his
%     voice (« the only one I could think of »), annotated in blue ink in his
%     hand. Equations numbered (1) to (24). Typed overstrikes that mask a
%     whole word are kept as \struck{}; single mistyped letters that the
%     typist overstruck in place are silently read as corrected.
%   Pages 14-18 — manuscript calculations for the theorem of the square and
%     the cube: n^*(L), the exact sequence for Pic of a product, a draft
%     definition of Pic_sp(X/Y), and the computation of L_{x_1…x_n} in terms
%     of the L_{x_i x_j}, L_{x_i}, L_∅, with the exponent α_n = (n-1)(n-2)/2.
%   Pages 19-20 — the beginning of his typed letter « Dear Murre »
%     (unramified functors, passage to the quotient, theorem of the cube,
%     separatedness of Corr_S(X,Y)); it runs on past page 20.
%
% Typescript conventions kept: underlined bold « Pic » (the functor) is
% \underline{\mathrm{Pic}}; the underlined O of the structure sheaf is
% \mathcal{O}; other underlined capitals (D, Z, L, J, G_m, A) are
% \underline{}. Words underlined by hand are \emph{}.
\documentclass[11pt,a4paper]{article}
\input{../preamble/grothendieck}

\folder{57}
\batch{1}
\pages{1}{20}
\dating{1962-[vers 1968]}
\watermark{Édition de démonstration}
\foldertitle{Picard : tapuscrit annoté (s.d.), lettres (s.d., 1962).}

\begin{document}

\section{Local Picard schemes}

\page{2}
\note{tapuscrit en anglais, annoté à l'encre bleue de sa main ; les numéros d'équation (1) à (3) sont ajoutés à la main}

Let $A$ be a complete local ring, $S = \mathrm{Spec}(A)$, $U = S - \{s\}$ the
complement of the closed point of $S$. We wish to interprete the Picard group
$\mathrm{Pic}(U)$ as the group $P(k)$ of $k$-valued points of a group-scheme
$P$ defined over the residue field $k$ of $A$. Assume for simplicity that $U$
is regular, and that there exists a resolution of singularities for $S$ not
changing $U$, i.e.\ a proper morphism
\[ f \colon X \longrightarrow S \]
inducing an isomorphism $f^{-1}(U) \simeq U$, such that $X$ be regular and
equal to the closure of $f^{-1}(U)$ (this open set will be identified with
$U$). Let $X_0$ be the special fiber of $X$, so that
\[ U = X - X_0 . \]
Because $X$ is regular, we have an exact sequence
\begin{equation*}
\tag{1} D \longrightarrow \mathrm{Pic}(X) \longrightarrow \mathrm{Pic}(U) \longrightarrow 0 ,
\end{equation*}
where $D$ is the group of divisors on $X$ concentrated on $X_0$ (which is a
free group of finite type $\underline{Z}^I$, $I$ being the set of irreducible
components of $X_0$). On the other hand, by the existence theorem for coherent
sheaves, it is known that
\begin{equation*}
\tag{2} \mathrm{Pic}(X) \simeq \varprojlim_n \mathrm{Pic}(X_n) ,
\end{equation*}
where $X_n = X \times_S S_n$, $S_n = \mathrm{Spec}(A/\underline{m}^{n+1})$.
Assume we can find group schemes $P_n$ over $k$ and isomorphisms
\begin{equation*}
\tag{3} \mathrm{Pic}(X_n) \xrightarrow{\ \sim\ } P_n(k) ,
\end{equation*}
and transition morphisms $P_m \to P_n$ ($m \geq n$) compatible with the
transition morphisms for the $\mathrm{Pic}(X_n)$. Denoting by $P$ the
pro-group $(P_n)$, or if one prefers, in case the transition morphisms
$P_m \to P_n$ are affine, the inverse limit (EGA IV 8) of the $P_n$, we get a
pro-group scheme resp.\ a group-scheme $P$, and an isomorphism

\page{3}
\begin{equation*}
\tag{4} \mathrm{Pic}(X) \xrightarrow{\ \sim\ } P(k) .
\end{equation*}
The map $D \to \mathrm{Pic}(X) \simeq P(k)$ défines a morphism of (pro) groups
\begin{equation*}
\tag{5} D_k \longrightarrow P ,
\end{equation*}
where $D_k$ is the constant group-scheme with value $D$. Assuming that the
cokernel of (5) exists in any reasonable sense a a pro-group scheme or a
group-scheme, let $Q$ be this kernel.\note{sic : « kernel » pour le conoyau qui vient d'être nommé} We get, in virtue of (1), an
injective homomorphism
\begin{equation*}
\tag{6} \mathrm{Pic}(U) \hookrightarrow Q(k) ,
\end{equation*}
the obstruction for an element of $Q(k)$ to be in the image being in
$H^1(k, D'_k)$, where $D'_k$ is the image of $D_k$ in $P$ (it is again a
constant group, with value a quotient $D'$ of $D$). \struck{If}
\add{If $k$ is sep.\ closed, or if} $D'$ is torsion-free, for example is $D'$
is equal to $D$ i.e.\ (5) is injective, then this obstruction vanishes, and
\struck{Q is a good candidate} (6) is an isomorphism. Thus $Q$ looks a good
candidate for a ``local Picard scheme''.

To construct the $P_n$'s, assume first that $A$ is a $k$-algebra, i.e.\ that
a lifting of \struck{the} \add{the} residue field $k$ exists, and has been
choosen. Then the natural candidate for $P_n$ (the only one I could think of,
indeed!) would be
\begin{equation*}
\tag{7} P_n = \underline{\mathrm{Pic}}_{X_n/k} .
\end{equation*}
We know from Murre that the right had side is indeed representable by a scheme
locally of finite type over $k$. There is however the usual trouble that we
have indeed a canonical \emph{injective} morphism
\begin{equation*}
\tag{8} \mathrm{Pic}(X_n) \hookrightarrow P_n(k) ,
\end{equation*}
but in general we are not sureit is bijective. However it is if $k$ is
separably closed; more specifically, the obstruction for an element of
$P_n(k)$ to be in the image of $\mathrm{Pic}(X_n)$ lies in the Brauer group
\struck{$\mathrm{Br}(k)$ of \ill{}}

\page{4}
of the Artin ring $B_n = H^0(X_n, \mathcal{O}_{X_n})$. Assuming $U$ hence $X$
hence $X_n$ to be connected, this Artin ring is local, and it's Brauer group
coincides with the Brauer group of its residue field $k_n$. Of course, for an
element of $P(k) = \lim P_n(k)$, these various obstructions match together and
come from an element of the Brauer group of the residue field
$k' = \bigcap_n k_n$ of the ring
$H^0(X, \mathcal{O}_X) = \varprojlim_n H^0(X_n, \mathcal{O}_{X_n})$,
\struck{of} which is just the normalisation of $A$. If for instance $A$ is
\emph{normal}, \add{then} $k' = k$ and the obstruction lies in
$\mathrm{Br}(k)$.

Whatever this be, defining $Q$ as before, we still do get an injective
homomorphism (6), but even if (5) is injective, we cannot be sure that (6) is
bijective. We are however if $k$ is separable closed. Therefore, we still can
more or less describe $Q$ (or at least $Q(k)$) in termes of local Picard
groups, in terms of \add{Galois} descent from the completion $A'$ of the
strict henselization of $A$: if $U' = S' - \{s'\}$, where
$S' = \mathrm{Spec}(A')$, and if
\begin{equation*}
\tag{9} \pi = \mathrm{Gal}(\bar{k}/k) ,
\end{equation*}
we will have an isomorphism
\begin{equation*}
\tag{10} Q(k) \simeq \mathrm{Pic}(U')^{\pi} ,
\end{equation*}
\struck{This is perfectly true provided $D'$ is torsion free, for instance (5)
injective.}\note{ligne tapée biffée d'un trait}
This is perfectly natural and in accordance with the familiar phenomena when
dealing with \emph{global} Picard schemes. Therefore $Q$ looks like the natural
object to be called a local Picard scheme.

To make sure it makes a sense, one will have however to analyze somewhat the
inverse system $(P_n)$, and the morphism (5). Moreover, to feel really secure,
one should check that up to canonical isomorphism, the group scheme $Q$
together with (10) does not depend on the choice of

\page{5}
the resolution; this should not be hard whenever we have existence of
resolutions in the strong sense of Hironaka\add{, so that two resolutions can
be dominated by a third one}, for instance in char.~0 of if $\dim A = 2$. I
should point out that, in any case, the algebraic structure thus defined on
$\mathrm{Pic}(U)$ depends in an essential way on the choice of the lifting of
the residue field; an instructive example on which to study the dependance of
this structure on the lifting is the one where $A$ is the completion of the
local ring at the origin of the projecting cone of an elliptic curve (here $Q$
is isomorphic to the elliptic curve in question, for the \emph{given} lifting
of $k$ …).

I have to correct an inaccuracy that has slipped in, when defining $Q$ in the
case when $k$ is not separably closed, so as to get (10). One will then have
to replace the constant group scheme $D_k$ by the twisted constant group
scheme $\underline{D}$, deduced by descent from the constant group-scheme
occuring after passage from $k$ to the separable closure.
\struck{In terms of the components $X_0^i$ of $X_0$, this can be defined as
the direct image, under the morphism
$\mathrm{Spec}(B_{0\,\mathrm{sep}}) \to \mathrm{Spec}(k)$, of the constant
sheaf $\underline{Z}$ on the spectrum of the largest separable subalgebra
$B_{0\,\mathrm{sep}}$ of $B_0$.}\note{passage encadré et barré à la main ;
« the components $X_0^i$ of $X_0$ » est tapé au-dessus d'une ligne biffée à la
machine, et la flèche de $\mathrm{Spec}(B_{0\,\mathrm{sep}}) \to
\mathrm{Spec}(k)$ manque au tapuscrit} This can be made explicit, in terms of
the separable algebraic closures $k_i$ of $k$ in the function fields of the
irreducible components $X_0^i$ of $X_0$, as the \struck{product of the
schemes} direct image of the constant sheaf $\underline{Z}$ under the
projection $\mathrm{Spec}(\prod k_i) \to \mathrm{Spec}(k)$.

To study the inverse system $(P_n)$, use the fact (SGA 6 XII) that the
transition morphisms are affine, which gives a sense to
$P = \varprojlim P_n$ as an actual group-scheme. It's Lie algebra is the
inverse limit of the Lie algebras of the $P_n$, which are
$H^1(X_n, \mathcal{O}_{X_n})$, therefore we get

\page{6}
\begin{equation*}
\tag{11} \underline{\mathrm{Lie}}(P) \simeq H^1(X, \mathcal{O}_X) .
\end{equation*}
As $R^1 f_*(\mathcal{O}_X)$ is coherent and concentrated at the origin, it
follows that its group of sections $H^1(X, \mathcal{O}_X)$ is of finite
dimension. In fact, it seems, with a little care, that replacing if necessary
$X$ by a suitable model dominating it, we may assume that there exists an
\add{invertible} ideal $J$ on $X$, defining a subscheme $Z$ of $X$ having same
support as $X_0$, and such that \add{$\underline{L} =$} $J/J^2$ be an
\emph{ample sheaf on $Z$}. This will imply that, replacing the system
$(X_n)_{n}$ by the system $(Z_n = V(J^{n+1}))_{n}$, which defines an
isomorphic pro-\struck{group} \add{object} of Picard-schemes, we obtain an
inverse system of group-schemes $P'_n = \underline{\mathrm{Pic}}_{Z_n/k}$
where for large $n$ the transition morphisms $P'_{n+1} \to P'_n$ are
\emph{isomorphisms} (because
$H^1(Z_0, \underline{L}^{\otimes n}) = H^2(Z_0, \underline{L}^{\otimes n}) = 0$).\note{la flèche de $P'_{n+1} \to P'_n$ manque au tapuscrit}
Thus it seems that, under the assumptions made \add{from the start}, one can
prove that the pro-group $(P_n)$ is essentially constant, hence its inverse
limit $P$ is in fact a \emph{group scheme locally of finite type over $k$}.
This is certainly so, at least, if $\dim A = 2$, $A$ normal, as follows from
the negative definiteness of the intersection matrix of the components of
$X_0$.

This $P$ seems to be definable in quite a reasonable way, and one still has to
investigate the morphism (5) (or more accurately, the morphisme
$\underline{D} \to P$ defined by descent from the case where $k$ is separably
closed; but then we may as well assume $k$ separably closed from the start,
\add{which we will do}). We would be particularly happy if this were a closed
immersion, or what amounts to the same, if the morphism
\begin{equation*}
\tag{12} D \longrightarrow \mathrm{NS}(X_0)
\end{equation*}
is injective, were
$\mathrm{NS}(X_0) = \underline{\mathrm{NS}}_{X_0/k}(k)$,
$\underline{\mathrm{NS}} = \underline{\mathrm{Pic}} / \underline{\mathrm{Pic}}^0$.
(The equivalence comes from the fact that
$P = \underline{\mathrm{Pic}}_{Z_n/k}$ for some large $n$, and that
\struck{Pic}

\page{7}
as recalled above, the morphism
$\underline{\mathrm{Pic}}_{Z_n/k} \to \underline{\mathrm{Pic}}_{X_0/k}$ is
affine, and hence $\mathrm{NS}(Z_n) \to \mathrm{NS}(X_0)$ has finite kernel.)
I expect that this is always the case, \struck{\ill{}} In case $A$ is normal
of dimension 2, this follows from Du-Val -- Mumford's negative definiteness.
It would he nice to have a proof in general (perhaps through a negative
definiteness of an intersection matric $X^i . C^j$, where for every
irreducible component $X^i$ of $X_0$, $C^i$ denotes a \add{suitable} curve on
$X^i$, for instance defined in terms of a suitable relatively ample sheaf on
$X$ by taking intersections of corresponding hypersurface sections on the
$X^i$'s ?). In case (12) is indeed injective, we have no trouble defining
\begin{equation*}
\tag{13} Q = P/\underline{D}
\end{equation*}
as a \emph{group-scheme locally of finite type over} $k$. We then have
\begin{equation*}
\tag{14} P^0 \xrightarrow{\ \sim\ } Q^0 ,
\end{equation*}
an isomorphism for the connected components, and
\begin{equation*}
\tag{15} Q/Q^0 = (P/P^0)/\underline{D} \simeq \underline{\mathrm{NS}}_{Z_n/k} / \underline{D} ,
\end{equation*}
where $n$ is chosen large enough. Of course (14) and (11) yield an isomorphism
\begin{equation*}
\tag{16} \underline{\mathrm{Lie}}(Q) \simeq H^1(X, \mathcal{O}_X) .
\end{equation*}
Let us remark by the way that we have an \emph{injective} restriction morphism
\begin{equation*}
\tag{17} H^1(X, \mathcal{O}_X) \hookrightarrow H^1(U, \mathcal{O}_U) = H^2_s(S)
\end{equation*}
as one sees by remarking that, with the above notations for $J$,
$\underline{L}$, \add{the sheaf} $i_*(\mathcal{O}_U)/\mathcal{O}_X$
($i \colon U \to X$ the inclusion) has a composition series whose quotients
are the $\underline{L}^{-n}$, $n \geq 1$, whose $H^0$ is zero, and the map
(17) is just the canonical map
$H^1(X, \mathcal{O}_X) \hookrightarrow H^1(X, i_*(\mathcal{O}_U))$.\note{l'astérisque de $i_*$ et la crosse de $\hookrightarrow$ sont ajoutés à la main}
Thus, the Lie algebra of the local Picard scheme is identified with a subspace
of $H^1(U, \mathcal{O}_U)$. (This suggests that if $A$ is Cohen-Macaulay of
dim.~$\geq 3$, $Q$ is discrete, hence $\mathrm{Pic}(U)$ countable etc.)

\page{8}
The previous constructions rely heavily on the existence and choice of a field
of representatives. Even this being granted, the question of caracterising $Q$
as an actual \emph{scheme} \add{over $k$}, i.e.\ as a functor on the category
of $k$-algebras, has not been touched. In fact, the description I can give of
this functor is \add{(tautological, and is)} in terms of a particular choice of
a resolution $X$, and it really gives \struck{rather} a description of $P$
rather than $Q$. One then still has to define $Q$ as the cokernel of the silly
morphism $\underline{D} \to P$. As a matter of fact, even if we restrict to
arguments $k'$ which are, say, finite extensions of $k$ (but not necessarily
separable ones), I cannot describe $Q(k')$ in terms of local Picard groups of
local rings as $A' = A \hat{\otimes}_k k'$. The trouble, of course, is that
$X \otimes_k k'$ will in general no longer be a resolution of $S'$, i.e.\ it
will no longer be regular. However, such geometric descriptions will be
possible provided we take as arguments algebras which are
\struck{smooth of formally smooth} \add{geometrically regular} over $k$, so
that tensoring with them will not destroy the regularity of $X$. (Thus, the
case when $k'$ is a power series ring, for instance the completion of the
local ring of $Q$ at the origin, $Q$ being smooth, is of interest …) Let's
now try to be specific. By defintion \add{more or less,} for any
\add{affine} scheme $T$ over $k$, we have
\begin{equation*}
\tag{18} P(T) = \varprojlim P_n(T) \hookleftarrow \varprojlim_n \mathrm{Pic}(X_{n_T}) \text{\add{$/\mathrm{Pic}(T)$}}
\qquad (X_{n\,T} = X_n \times_k T) .
\end{equation*}
\marginal{affineness of $T$ used to \uncertain{have}
$H^1(T, \mathcal{O}_T) = H^2(T, \mathcal{O}_T) = 0$}
In fact, it is not hard to check that the left hand side functor is the étale
sheaf (on $(\mathrm{Sch})_{/k}$) associated to the right-hand side functor.
This comes from the more precise statement that the obstruction, for an
element of $P(T)$, to come from an element of
$\varprojlim_n \mathrm{Pic}(X_{nT})$, lies in
$H^2(T_{\text{ét}}, \underline{G}_m)$ (at least for $T$ affine, which is used
to insure that it's $H^1$ and $H^2$ with values in $\mathcal{O}_T$ vanish.
For simplicity, I have assumed here that $A$ is normal). As we have
seen\note{dans (18), « (18) », la flèche $\hookleftarrow$, le $n$ sous la
seconde limite et « $/\mathrm{Pic}(T)$ » sont de sa main, ce dernier
remplaçant une fin de formule biffée ; dans la parenthèse finale, les
exposants 2 et 3 tapés sont corrigés à la main en 1 et 2}

\page{9}
on the other hand that $P$ is locally of finite presentation over $k$, it
commutes (as a functor) to filtering direct limits of rings, and therefore it
is known when we know it's restriction to arguments which are of finite type
over $k$, and a fortiori noetherian. For such \add{(i.e.\ noetherian)}
arguments, using the existence theorem for coherent sheaves, we can give a
more geometric expression for the right hand side of (18), \struck{namely}
Assume for simplicity $T = \mathrm{Spec}(B)$ affine, and let
\[ C = A \hat{\otimes}_k B = \lim A_n \otimes_k B , \]
which is an adic noetherian ring augmented to $B$ (the augmentation ideal
being an ideal of definition). Consider \struck{\ill{}}
\add{$X \hat{\otimes}_k B$} $= X \otimes_A C$, which is a proper scheme over
$C$. We then have
\begin{equation*}
\tag{19} \mathrm{Pic}(X \hat{\otimes}_k B) \simeq \varprojlim_n \mathrm{Pic}(X_{nB}) ,
\end{equation*}
\struck{\ill{}} \add{$\mathrm{Pic}(X \hat{\otimes}_k B)/\mathrm{Pic}(B) \hookrightarrow P(B)$}
and thus on \emph{noetherian} arguments $B$, $P$ appears as the étale sheaf
associated to the functor
$B \mapsto \mathrm{Pic}(X \hat{\otimes}_k B)$\add{$/\mathrm{Pic}(B)$},
\add{which is a separated presheaf,} this statement being complemented by the
description given above of the group for obstructions for an element of $P(B)$
to belong to $\mathrm{Pic}(X \hat{\otimes}_k B)$.

As for $B \mapsto Q(B)$, this appears as the étale sheaf associated with the
functor
\begin{equation*}
\tag{20} B \longmapsto \mathrm{Pic}(X \hat{\otimes}_k B)/\bigl(\text{\add{$\mathrm{Pic}(B) +$}}\, \underline{D}(B)\bigr) \quad (\hookrightarrow Q(B))
\end{equation*}
where $\underline{D}$ is the twisted constant group over $k$ defined above.
For an element of $Q(B)$ to belong to the first hand side there are two
successive obstructions, the first is in $H^1(\mathrm{Spec}(B), \underline{D})$,
which vanishes \add{for instance} if $\underline{D}$ is constant (for instance
$k$ separably closed) and $B$ normal, the second is in the Brauer group of
$B$, and vanishes \add{for instance} if $X_0$ has a zero-cycle of degree~1.

Assume for simplicity that the irreducible components of $X_0$ are

\page{10}
geometrically irreducible, i.e.\ $\underline{D}$ is constant, and assume $B$
geometrically regular over $k$. This second condition implies that
$X \hat{\otimes}_k B$ is regular, and denoting by $U \hat{\otimes}_k B$ the
inverse image of $U$ in this scheme, \struck{the first condition} this
together with the first condition implies that
\struck{the left hand side of (20)}
\add{$\mathrm{Pic}(X \hat{\otimes}_k B)/D(B)$} is isomorphic to
$\mathrm{Pic}(U \hat{\otimes}_k B)$. This shows that on the sub-category of
arguments $B$ which are noetherian and geometrically regular, $Q$ is the étale
sheaf associated to the presheaf
\begin{equation*}
\tag{21} B \longmapsto \mathrm{Pic}(U \hat{\otimes}_k B)\text{\add{$/\mathrm{Pic}(B)$}} \quad (\hookrightarrow Q(B)) .
\end{equation*}
The obstruction for an element of $Q(B)$ to belong to
$\mathrm{Pic}(U \hat{\otimes}_k B)$\add{$/\mathrm{Pic}(B)$} is in the Brauer
group $\mathrm{Br}(B)$. Thus, under very stringent conditions on $B$ at least
(being geometrically regular) we get a description of $Q(B)$ in terms of
actual local Picard groups. I should have stated that $U \hat{\otimes}_k B$
can also be interpreted as the inverse image of $U$ in the scheme
$\mathrm{Spec}(A \hat{\otimes}_k B)$, and thus makes a sense independently of
the choice of a particular resolution.

One may wish to construct a local Picard scheme \struck{without} also in case
$A$ is not equi-caracteristic, which implies that $\mathrm{car}.\,k > 0$. It
seems likely that this can be done if $k$ is perfect. The key-point here is
the construction of schemes $P_n$, which leads us to the following

\emph{Problem}. Let $A$ be a local Artin ring with perfect residue field $k$
of car.~$p > 0$, $X$ a proper scheme over $A$. Give a ``natural''
construction of a group scheme locally of finite type $P$ over $k$, together
with an imbedding
\begin{equation*}
\tag{22} \mathrm{Pic}(X) \hookrightarrow P(k) ,
\end{equation*}
the obstruction for an element of $P(k)$ to beling to $\mathrm{Pic}(X)$
\add{(}being in $\mathrm{Br}\, H^0(X, \mathcal{O}_X)$, and hence\add{)}
vanishing if $k$ is algebraically closed.\note{l'énoncé du problème est marqué
d'un trait vertical dans la marge de gauche}

\page{11}
Here is a candidate for a functor, which may turn out to be representable, in
which case this would be $P$. Let $\underline{A}$ be the ring-scheme over $k$
defined by $A$ (see Greenberg's papers, or Serre's in Bulletin Soc.\ Math.),
so that
\begin{equation*}
\tag{23} A \simeq \underline{A}(k) .
\end{equation*}
For every algebra $B$ over $k$, consider the ring $\underline{A}(B)$, which is
an algebra over $\underline{A}(k) = A$. We may thus consider
\[ X_B \overset{\mathrm{dfn}}{=} X \otimes_A \underline{A}(B) , \]
and consider the functor
\begin{equation*}
\tag{24} B \longmapsto \mathrm{Pic}(X_B) .
\end{equation*}
\note{le numéro tapé « (34) » est corrigé à la main en « (24) »}
This of course may not be representable, even if $A = k$. However, it may turn
out that the fppf sheaf associated to this functor is representable, and this
then would give the looked for candidate for $P$. In any case, Artin's work
gives us a very handy set of necessary and sufficient conditions for a group
functor to be representable by a scheme locally of finite type, and in
principle it should be possible to decide wether or not the
\struck{functor} sheaf just considered satisfies to these conditions or not.
If so, one may expect that the developments on local Picard schemes given for
the case when there is a field of representatives will carry over to the case
when the residue field is perfect with car.~$> 0$.

\page{12}
In the construction of a local Picard scheme $Q$, we have seemed to use in a
technically essential way the \struck{fact} \add{assumption} that $U$ is
regular (plus resolution). I would conjecture that a reasonable local Picard
scheme $Q$, \add{\emph{locally of finite type over $k$},} can be constructed
also without this assumption (provided that a field of representatives is
given, or the residue field is perfect of char.~$p > 0$, of course). Instead
of a resolution of singularities, one may then wish to find a proper
surjective morphism
\[ f \colon X \longrightarrow S \]
inducing an isomorphism $f^{-1}(U) \xrightarrow{\ \sim\ } U$, and such that
the corresponding map
\[ \mathrm{Pic}(X) \longrightarrow \mathrm{Pic}(U) \]
is \emph{surjective}. We then could repeat the construction of $P$, as
$\lim \underline{\mathrm{Pic}}_{X_n/k}$, as above, and of $Q$ as
$\mathrm{Coker}(\underline{D} \to P)$, where $\underline{D}$ is a group
scheme corresponding to divisors concentrated on $X_0$ (possibly no longer
discrete ?). Another way of approach would be to take a resolution of
singularities $f \colon X \to S$, but where now we cannot assume any longer
that $f^{-1}(U) \to U$ is an isomorphism, and to describe a scheme $P$ in
terms of invertible sheaves on $X$, together with \emph{descent data} of
$\underline{L}|f^{-1}(U)$ to $U$. A nice test for the yoga of local Picard
schemes without regularity assumptions on $U$ would be to see if it is true
in general that for $A$ Cohen-Macaulay of $\dim \geq 3$ (more generally, for
$H^2_{\underline{m}}(A) = 0$), the group $\mathrm{Pic}(U)$ is indeed
countable, and for separably closed residue field \struck{\ill{}}
\add{is invariant under} \emph{separable} field extension $k \to k'$, $A$
being replaced by $A \hat{\otimes}_k k'$ (which would express the expected
discrete structure of the local Picard scheme).

\section{Carré, cube, rigidité…}
\note{titre de sa main, en haut à droite d'une feuille de garde par ailleurs
blanche (p.~13), suivi de deux croix}

\page{14}
\note{feuillet de calculs, encres noire et bleue}
\[
\boxed{\, n^*(L) \simeq L^{\otimes n}\, (\delta L)^{\frac{n(n-1)}{2}} \,}
\qquad\qquad
n + 2\,\frac{n(n-1)}{2} = n^2
\]
\[
\delta L = L^{\otimes 2}\, \delta'(L) \qquad \delta' L_x = L_{2x}\, L_x^{-4}
\]
\[
\boxed{\, n^*(L) \simeq L^{n^2}\, (\delta' L)^{\frac{n(n-1)}{2}} \,}
\]
\note{l'exposant de $L$ dans la seconde formule encadrée est noyé sous une
surcharge ; on y lit $n^2$ sous réserve}
\[ A \longrightarrow A \times A \qquad A \longrightarrow \hat{A} \]
\[
\boxed{\, \underline{\underline{\mathrm{Pic}}}^{\tau}_{(\prod X_i)/S} \simeq \prod_i \underline{\underline{\mathrm{Pic}}}^{\tau}_{X_i/S} \,}
\]
\add{sans torsion}\note{ajouté en travers, souligné deux fois, à côté de l'encadré}
\[
0 \to \prod_i \underline{\underline{\mathrm{Pic}}}_{X_i/S} \longrightarrow \underline{\underline{\mathrm{Pic}}}_{(\prod X_i)/S} \longrightarrow \prod_{\{i,j\}} \mathcal{C}_S(X_i, X_j) \longrightarrow 0
\]
\marginal{$\{i,j\}$ parcourant les parties à deux éléments de l'ensemble
d'indices $I$}\note{la note est écrite sous le produit, dans une bulle qui
l'entoure}

[\emph{\uncertain{Scindée}} si les $X_i/S$ ont tous une \emph{section}]

Suite exacte de faisceaux [pour la \ill{} \uncertain{essent.} fid.\ plate
qu.-compacte]

Les $X_i$ \uncertain{sont} propres \add{et plats}, $f_*(\mathcal{O}_{X_i})
\simeq \mathcal{O}_S$ universellement.

\marginal{\uncertain{explications} \ill{}}

\page{15}
\note{feuillet surchargé : à gauche, un schéma vertical $X \to Y$ et une
lettre $G$ noircie ; sur la droite, des traits de crayon en diagonale et un
encadrement ondulé traversent le texte}

$\mathrm{Pic}(X/Y)$

$\mathrm{Pic}_{\mathrm{sp}}(X/Y)$ : classes à \uncertain{isom.\ près} de
faisceaux inversibles $\underline{L}$, \uncertain{sections}
\uncertain{marquées}, telles que \struck{\ill{} soit} (i) l'on ait
$H^i_f(\underline{L}) = 0$ \uncertain{pour} $i > 0$, \struck{\ill{}}

donc $\underline{H}^0_f(\underline{L})$ est loc.\ libre de rang
$\chi_{X/Y}(\underline{L})$. (ii) \struck{$\chi_{X/Y}(\underline{L})$}
\ill{} $\mapsto 1$, i.e.\ l'\ill{} \ill{} équivaut à 2 sections
\uncertain{marquées} \ill{} $H^0_f(\underline{L}) \to$ \ill{}

$\underline{H}^0_f(\mathcal{O}_X) = \mathcal{O}_Y$ « universellement »

\ill{} \uncertain{sections marquées} \ill{}

\drawing{un diagramme à demi biffé, en partie écrit tête-bêche : des produits
$P^1 \times P$, $P^0 \times_P P \rightrightarrows P_1 \to P$, des flèches
verticales et une flèche courbe ; les termes ne se lisent pas assez pour être
composés}

Une telle donnée équivaut \uncertain{à} \ill{} : la donnée \ill{}

\page{16}
\[
\begin{array}{ccc}
 & X_i \times X_j & \longrightarrow\ \prod X_i \\
 \nearrow & & \nearrow \\
X_{[ijkl]} \longrightarrow & X_k \times X_l &
\end{array}
\]
\note{diagramme recomposé : $X_{[ijkl]}$ (sur un $X$ surchargé) s'envoie dans
$X_i \times X_j$ et $X_k \times X_l$, qui s'envoient dans $\prod X_i$}
\[
X_{ijkl} = \begin{cases}
X_{\emptyset} & \text{si } \{i,j\} \cap \{k,l\} = \emptyset \\
X_{\alpha} & \text{si } \{i,j\} \cap \{k,l\} = \{\alpha\}
\end{cases}
\]
\[ L_{m_1 x_1 + \cdots + m_n x_n} \]
\[ L_{x,y} \simeq L_{xe}\, L_{ye}\, \mathcal{M}_{xy} \qquad \mathcal{M}_{xy}\, \mathcal{M}_{yx} \]
\note{au-dessus de $L_{x,y}$, $L_{xe}$, $L_{ye}$, trois petits exposants
surchargés, illisibles}

\struck{$L_{xy}\, L_{yz}$} \qquad \struck{$L_{x+y+z}$}
\[
L_{x_1 + \cdots + x_n} = \Bigl(\prod_{i<j} L_{x_i + x_j}\Bigr) \Bigl(\prod L_{x_i}^{-n+2}\Bigr) L_{\emptyset}^{\alpha_n}
\]
\[
L_{x+y} = L_x\, L_y\, (\delta L)_{x,y}
\]
\marginal{symétrique}\note{écrit près de $(\delta L)_{x,y}$, entouré}
\[
\prod_{i \neq j} L_{x_i + x_j} = \Bigl(\prod L_{x_i}^{n-1}\Bigr) \prod_{i \neq j} (\delta L)_{x_i, x_j}
\]
\[
\boxed{\, L_{x_1 + \cdots + x_n} = \prod_{i<j} (\delta L)_{x_i, x_j} \prod_i L_{x_i}\, L_{\emptyset}^{\alpha_n} \,}
\]
\note{dans l'encadré, un facteur surchargé avant le premier produit et un
autre, biffé, avant $\prod_i L_{x_i}$ ; sous l'encadré, une accolade et
l'annotation « si $x_i = 0$ pour \ill{} » mènent au calcul suivant}
\[
L_{\emptyset}^{-\frac{n(n-1)}{2}}\, L_{\emptyset}^{\frac{(n-1)(n-2)}{2}} = L_{\emptyset}^{-n+1}
\qquad (\delta L)_{0,0} = L_{\emptyset}^{-1}
\]
\[
1 = -\frac{n(n-1)}{2} + n + \frac{(n-1)(n-2)}{2}
\]
\note{en marge droite, deux autres essais de cette identité, biffés}

\page{17}
\[ 1 + 1 - 1 = 1 \qquad -(n-2) - (n-2) + (n-3) = -(n-1) \]
\[
\alpha_n + \alpha_n - \alpha_{n-1} + 1 = \text{\struck{\ill{}}}\quad 2\alpha_n - \alpha_{n-1} + 1 = \alpha_{n+1}
\]
\[
\alpha_3 = 1 \quad \alpha_4 = 3 \quad \alpha_5 = 6 \quad \alpha_6 = 10 \quad \alpha_7 = 15
\]
\struck{$\alpha_n = an + b$, $2\alpha_n - \alpha_{n-1} + 1$ $= 2an + 2b - an + a - b + 1$ $= an + (\dots)$, $\alpha_{n+1} = an + a + b$}\note{essai linéaire barré de deux grandes croix ; la parenthèse du terme constant est surchargée}
\begin{align*}
\alpha_n &= an^2 + bn + c \\
2\alpha_n - \alpha_{n-1} + 1 &= 2an^2 + 2bn + 2c - a(n^2 - 2n + 1) - b(n-1) - c + 1 \\
&= an^2 + (b + 2a)n + (c - a + b + 1) \\
\alpha_{n+1} &= a(n^2 + 2n + 1) + b(n+1) + c \\
&= an^2 + (2a + b)n + (a + b + c)
\end{align*}
\[
a + b + c = -a + b + c + 1 \qquad a = -a + 1 \qquad a = \tfrac12
\]
\[
\alpha_n = \tfrac12 n^2 + bn + c \qquad \alpha_2 = 0,\ \alpha_3 = 1
\]
\[
\tfrac12 \cdot 4 + 2b + c = 0 \qquad \tfrac12 \cdot 9 + 3b + c = 1
\]
\[
2b + c = -2 \qquad 6b + 2c = -7 \qquad 2b = -3 \qquad c = +1
\]
\note{plusieurs chiffres de ce calcul sont surchargés ($\alpha_3$ sur
$\alpha_2$, $-7$, $-3$, $+1$, $\pm 1$ dans $c - a + b + 1$) ; on donne la
dernière lecture}
\[
\tfrac12 n^2 - \tfrac32 n + 1 = \tfrac12 [n^2 - 3n + 2]
\qquad
\boxed{\, \alpha_n = \frac{(n-2)(n-1)}{2} \,}
\]
\struck{$L_{x_1 \dots x_n}$ $\simeq \bigl(\prod_{i \neq j} L_{x_i x_j}\bigr) \bigl(\prod_i L_{x_i}^{-1}\bigr) L_{\emptyset}^{\frac{(n-1)(n-2)}{2}}$}
\[
\boxed{\, L_{x_1 \dots x_n} \simeq \Bigl(\prod_{i \neq j} L_{x_i x_j}\Bigr) \Bigl(\prod_i L_{x_i}\Bigr)^{-n+2} L_{\emptyset}^{\frac{(n-1)(n-2)}{2}} \,}
\]

\page{18}
\[
\begin{matrix} x & y & z \\ a & b & c \end{matrix}
\qquad
\boxed{\, L_{x,y,z} \simeq L_{xy}\, L_{y,z}\, L_{x,z}\, L_x^{-1} L_y^{-1} L_z^{-1} L_{\emptyset} \,}
\]
\marginal{« Th.\ du cube » ---}
\begin{align*}
L_{xyzt} &= L_{(x,y)z}\, L_{z,t}\, L_{(x,y),t}\, L_{(x,y)}^{-1} L_z^{-1} L_t^{-1} L_{\emptyset} \\
&= L_{xy} L_{y,z} L_{xz} L_x^{-1} L_y^{-1} L_z^{-1} L_{\emptyset}\, L_{zt} \\
&\qquad L_{xy} L_{y,t} L_{x,t} L_x^{-1} L_y^{-1} L_t^{-1} L_{\emptyset}\,
L_{xy}^{-1} L_z^{-1} L_t^{-1} L_{\emptyset}
\end{align*}
\note{les facteurs sont cochés un à un et ceux qui se compensent biffés ;
on n'a pas reproduit ces marques}
\[
\boxed{\, L_{xyzt} = L_{xy} L_{xz} L_{xt} L_{yz} L_{yt} L_{zt}\,
L_x^{-2} L_y^{-2} L_z^{-2} L_t^{-2}\, L_{\emptyset}^{3} \,}
\]
\[
L_{xyztu} = L_{xyzt}\, L_{tu}\, L_{xyzu}\, L_{xyz}^{-1} L_t^{-1} L_u^{-1} L_{\emptyset}
\]
\note{suivent trois lignes de développement, cochées et biffées de même}
\[
\boxed{\,
\begin{array}{l}
L_{xyztu} = L_{xy} L_{xz} L_{xt} L_{xu} L_{yz} L_{yt} L_{yu} L_{zt} L_{zu} L_{tu} \\
\qquad L_x^{-3} L_y^{-3} L_z^{-3} L_t^{-3} L_u^{-3}\, L_{\emptyset}^{6}
\end{array}\,}
\]
\[
L_{x_1 \dots x_n} = \prod_{i \neq j} L_{x_i x_j} \prod L_{x_i}^{-(n-2)} L_{\emptyset}^{\alpha_n}
\]
\begin{align*}
L_{x_1 \dots x_{n+1}} &= L_{(x_1 \dots x_{n-1}) x_n x_{n+1}} \\
&= L_{x_1 \dots x_n}\, L_{x_1 \dots x_{n-1} x_{n+1}}\, L_{x_n x_{n+1}}\,
L_{x_1 \dots x_{n-1}}^{-1} L_{x_n}^{-1} L_{x_{n+1}}^{-1} L_{\emptyset} \\
&= \prod_{1 \leq i < j \leq n} L_{x_i x_j} \prod_{1 \leq i \leq n} L_{x_i}^{-(n-2)} L_{\emptyset}^{\alpha_n}
\prod_{\substack{1 \leq i < j \leq n+1 \\ i, j \neq n}} L_{x_i x_j}
\prod_{\substack{1 \leq i \leq n+1 \\ i \neq n}} L_{x_i}^{-(n-2)} L_{\emptyset}^{\alpha_n} \\
&\qquad L_{x_n x_{n+1}} \prod_{1 \leq i < j \leq n-1} L_{x_i x_j}^{-1}
\prod_{1 \leq i \leq n-1} L_{x_i}^{+(n-3)} L_{\emptyset}^{-\alpha_{n-1}}\,
L_{x_n}^{-1} L_{x_{n+1}}^{-1} L_{\emptyset}
\end{align*}
\note{les bornes des produits sont lues sous réserve ; un « $1 \leq i, j
\leq n-1$, $i \neq j$ » est écrit à part, à gauche}
\[
\Bigl(\prod_{i \neq j} L_{x_i x_j}\Bigr) \Bigl(\prod L_{x_i}^{-(n-1)}\Bigr) L_{\emptyset}
\]
\note{la dernière ligne s'arrête là ; sous $\prod$, « $i \neq j$ » surcharge
un autre indice}

\section{Lettre à Murre}

\page{19}
\note{lettre tapée, non datée sur ces deux pages ; elle se poursuit au-delà
de la p.~20, hors de ce lot}

Dear Murre,

I am glad to hear that you are still willing to give the talk on unramified
functors. Here what I can say to your questions.

1. The theorem about passage to quotient I alluded to is the following:

\emph{Theorem}. Let $f \colon X \to Y$ be a morphisme of $S$-preschemes,
assume either $X$ and $Y$ \struck{of} \add{locally} finite presentation over
\struck{$Y$} $S$\note{lecture du $S$ surchargé sous réserve}, or $Y$ loc noeth
and $X$ \add{locally} of finite type over $Y$. Assume that the equivalence
relation \add{$R =$} $X \times_Y X$ defined by $f$ is flat over $X$ i.e.\
$\mathrm{pr}_1 \colon X \times_Y X \to X$ is flat. Then the quotient $X/R$
exists in the strongest reasonable sense, i.e.\ one can factor $f$ into a
compositum $X \to Z \to Y$, with $X \to Z$ faithfully flat locally of finite
presentation, $Z$ locally of finite \struck{type} \add{pres.} over
\struck{$S$} $Y$\note{lecture sous réserve de la lettre surchargée} (in fact
of finite presentation over $S$ if $X$ is so) and $Z \to Y$ a
monomorphism.\note{les flèches de $X \to Z \to Y$ et de $X \to Z$ manquent au
tapuscrit}

Of course the factorization is unique, andthe theorem can be expressed by
saying that the quotient sheaf (for the fpqc topology) $X/R$ is
representable. That is in fact how the theorem is proved.

Raynaud has recently made a very nice (and non trivial) application of this
theorem, by proving the following: if $S$ is the spectrum of a discrete
valuation ring, $G$ a group prescheme of finite type over $S$, $H$ a closed
\add{and flat} sub-group scheme, such that $G_t/H_t$ is \struck{affine}
quasi-affine (where $t$ is the generic point of $S$) then $G/H$ is
representable as a quasi-affine and flat $S$-scheme, which is even affine if
$H$ is invariant (i.e.\ if $G$ is a flat group scheme of finite type with
affine generic fibre, than $G$ is affine). This extends immediately to a base
which is regular of dim one. Raynaud is now trying to extend his construction
to the case when he drops the quasi-affinness assumption, namely to construct
still $G/H$

\page{20}
as a quasi-projective scheme over $S$.

2. Theorem of the cube. I believe we discussed about it time ago, but maybe
the proof I told you was valid only if one assumes the Pic functor of one of
the \struck{\ill{}} \add{factors} involved representable. To prove
unramifiedness of the functor $\underline{\mathrm{Corr}}$ however you need
only a week infinitesimal form of the theorem of the square, for which you
will find a proof in the manuscript notes I am joining on correspondance
classes, containing also the proofs of the statements you were recalling in
your question 4. I hope you will be able to read them, I agree the
handwriting is wretched and the notes moreover very sketchy. -- On the other
hand, I recall you that the theorem of the cube follows rather formally once
one knows separatedness of $\underline{\mathrm{Corr}}_S(X,Y)$ for two of the
three factors involved, and using the usual formal properties of the Picard
functor (among which commutation with inverse limits of Artin rings is the
less trivial).

3. As for the separatedness of $\underline{\mathrm{Corr}}_S(X,Y)$, this is
about trivial whenever the Pic functor of one of the factors $X,Y$ is
separated? Now this is certainly the case for $X$ if its geometric fibers are
integral, (a fortiori if $X$ is an abelian scheme over $S$ !). To show this,
one may assume $S$ the spectrum of a \add{discrete} valuation ring, and one is
reduced to show that if $\underline{L}$ is an invertible sheaf on $X$ whose
restriction to the general fiber $X_t$ is trivial, then $\underline{L}$ is
trivial. Now $X_1$ is an open subset of $X$, and the assumption on
$\underline{L}$ can be expressed by saying that $\underline{L}$ is defined by
a Cartier divisor whose support is contained in the special fiber $X_0$. Now
$X_0$ itself is already a Cartier divisor (defined by a global equation
$t = 0$) and moreover is an integral subscheme of $X$, from this follows that
the divisor $D$ is a multiple of $X_0$ (assume for
\note{la phrase se poursuit sur la page suivante, hors de ce lot}

\end{document}
